Пусть $\varphi : F \to F$ -- автоморфизм поля $F$ (изоморфизм поля на себя).

а) Пусть $char\ F = 0$. Верно ли, что $\varphi$ сохраняет $\Q$? 
(то есть при $q \in \Q$ выполнено равенство $\varphi(q) = q$). 

Да, это верно: $\varphi(1) = 1 \Rightarrow \forall n \in \N \varphi(n) = n$.
Так как $\varphi(a) = \varphi(a)^{-1}$, то понятно, что

$$
\varphi \left( \frac{m}{n} \right) = \varphi(m) \cdot \varphi(n)^{-1} = 
m \cdot n^{-1} = \frac{m}{n}
$$

б) Пусть $char\ F = p$. Верно ли, что $\varphi$ сохраняет $\Z_p$?

Аналогично пункту 1:

$$
\forall a \in \Z_p \ \varphi(a) = a \cdot \varphi(1) = a
$$

Пусть теперь $F \supset \Q$.

в) Если $\alpha \in F$ -- корень многочлена $f(x) \in \Q[x]$, 
то $\varphi(\alpha)$ также корень многочлена $f(x)$.

Пусть $f(x) = a_n x^n + \dots + a_1 x + a_0$. Применим $\varphi$ к $f(x)$:

$$
\varphi(a_n x^n + \dots + a_1 x + a_0) = 
\varphi(a_n) \varphi(x^n) + \dots + 
\varphi(a_1) \cdot \varphi(x) + \varphi(a_0) = 
a_n (\varphi(x))^n + \dots + a_1 \cdot \varphi(x) + a_0
$$

Все коэффициенты из $\Q$, поэтому по пункту а) $\varphi$ их не изменил.
Отсюда видно, что $\varphi(\alpha)$ также является корнем $f(x)$.

г) Алгебраические элементы при автоморфизмах переходят в сопряжённые 
(над $\Q$) элементы.

Возьмём алгебраический элемент $\alpha \in F$. Обозначим его минимальный
многочлен за $f(x)$, $f(x) \in \Q[x]$. Из пункта в) следует, что
$\varphi(\alpha)$ также является корнем этого многочлена.
Так как $\varphi(\alpha)$ является корнем минимального многочлена элемента
$\alpha$, то их минимальные многочлены совпадают. По определению это и
означает, что $\alpha$ и $\varphi(\alpha)$ сопряжены.
